\input{preamble.tex}

\title{Лабораторная работа No. 12. \\Настройка NAT}
\author{Сущенко Алина Николаевна\\ НПИбд-01-23}
\institute{Российский университет дружбы народов имени Патриса Лумумбы}
\date{2026}

\begin{document}

\frame{\titlepage}

\begin{frame}
  \frametitle{Цель работы}
  \begin{itemize}
  \item Настроить преобразование сетевых адресов для обеспечения доступа локальных сетей организации во внешнюю сеть (Интернет) и организовать доступ из внешней сети к внутренним серверам (WWW, почтовый, файловый, RDP).
  \end{itemize}
\end{frame}

\begin{frame}
  \frametitle{Первоначальная настройка маршрутизатора ansusenko-gw-1}
  \centering
  \includegraphics[width=\textwidth]{../images/01.png}
  \captionof{figure}{Первоначальная настройка маршрутизатора ansusenko-gw-1.}
\end{frame}

\begin{frame}
  \frametitle{Первоначальная настройка коммутатора ansusenko-sw-1}
  \centering
  \includegraphics[width=\textwidth]{../images/02.png}
  \captionof{figure}{Первоначальная настройка коммутатора ansusenko-sw-1.}
\end{frame}

\begin{frame}
  \frametitle{Настройка интерфейсов маршрутизатора ansusenko-gw-1}
  \centering
  \includegraphics[width=\textwidth]{../images/03.png}
  \captionof{figure}{Настройка интерфейсов маршрутизатора ansusenko-gw-1.}
\end{frame}

\begin{frame}
  \frametitle{Настройка интерфейсов коммутатора ansusenko-sw-1}
  \centering
  \includegraphics[width=\textwidth]{../images/04.png}
  \captionof{figure}{Настройка интерфейсов коммутатора ansusenko-sw-1}
\end{frame}

\begin{frame}
\frametitle{Настройка интерфейсов маршрутизатора msk-ansusenko-gw-1}
    \centering
    \includegraphics[width=\textwidth]{../images/05.png}
    \captionof{figure}{Настройка интерфейсов маршрутизатора msk-ansusenko-gw-1 и маршрута по умолчанию.}
\end{frame}

\begin{frame}
\frametitle{Настройка пула адресов}
    \centering
    \includegraphics[width=\textwidth]{../images/06.png}
    \captionof{figure}{Настройка пула адресов.}
\end{frame}

\begin{frame}
\frametitle{Настройка списка доступа}
    \centering
    \includegraphics[width=\textwidth]{../images/07.png}
    \captionof{figure}{Настройка списка доступа.}
\end{frame}

\begin{frame}
\frametitle{Правила доступа для сети "Донская"}
\centering
    \includegraphics[width=\textwidth]{../images/08.png}
    \captionof{figure}{Правила доступа для сети "Донская".}
\end{frame}

\begin{frame}
\frametitle{Доступ для сети кафедры}
    \centering
    \includegraphics[width=\textwidth]{../images/09.png}
    \captionof{figure}{Правила доступа для сети кафедры.}
\end{frame}

\begin{frame}
\frametitle{Доступ для сети администрации}
    \centering
    \includegraphics[width=\textwidth]{../images/10.png}
    \captionof{figure}{Правила доступа для сети администрации}
\end{frame}

\begin{frame}
\frametitle{Доступ для компьютера администратора}
    \centering
    \includegraphics[width=\textwidth]{../images/11.png}
    \captionof{figure}{Правила доступа для компьютера администратора.}
\end{frame}

\begin{frame}
\frametitle{Настройка PAT и интерфейсов}
 \centering
    \includegraphics[width=\textwidth]{../images/12.png}
    \captionof{figure}{Настройка PAT и интерфейсов.}
\end{frame}

\begin{frame}
\frametitle{Статическое преобразование для WWW-сервера}
    \centering
    \includegraphics[width=\textwidth]{../images/13.png}
    \captionof{figure}{Статическое преобразование для WWW-сервера.}
\end{frame}

\begin{frame}
\frametitle{Статическое преобразование для файлового сервера}
    \centering
    \includegraphics[width=\textwidth]{../images/14.png}
    \captionof{figure}{Статическое преобразование для файлового сервера.}
\end{frame}

\begin{frame}
\frametitle{Статическое преобразование для почтового сервера}
    \centering
    \includegraphics[width=\textwidth]{../images/15.png}
    \captionof{figure}{Статическое преобразование для почтового сервера.}
\end{frame}

\begin{frame}
\frametitle{Статическое преобразование для RDP-доступа}
    \centering
    \includegraphics[width=\textwidth]{../images/16.png}
    \captionof{figure}{Статическое преобразование для RDP-доступа к компьютеру администратора.}
\end{frame}

\begin{frame}
\frametitle{Контрольные вопросы}

\begin{enumerate}
    \item \textbf{В чём состоит основной принцип работы рассматриваемой технологии?}
    
    Основной принцип заключается в замене частных IP-адресов на публичные при выходе во внешнюю сеть и обратной замене при возврате ответов.
    
    \item \textbf{На каком оборудовании выполняется настройка?}
    
    Настройка выполняется на пограничном маршрутизаторе.
    
    \item \textbf{Можно ли применить средства Cisco IOS к субинтерфейсам?}
    
    Да, можно. Команды настройки применяются непосредственно к субинтерфейсам.
    
    \item \textbf{Что такое пулы IP-адресов?}
    
    Пул IP-адресов -- это диапазон публичных адресов, используемый для динамической трансляции частных адресов.
    
    \item \textbf{Что такое статические преобразования?}
    
    Статические преобразования -- это фиксированное соответствие между частным и публичным IP-адресом.
\end{enumerate}

\end{frame}

\begin{frame}
\frametitle{Вывод}
В ходе лабораторной работы удалось настроить преобразование сетевых адресов для обеспечения доступа локальных сетей организации во внешнюю сеть (Интернет) и организовать доступ из внешней сети к внутренним серверам (WWW, почтовый, файловый, RDP).
\end{frame}

\end{document}
